\documentclass[11pt,a4paper]{article}

\usepackage[margin=2.5cm]{geometry}
\usepackage[T1]{fontenc}
\usepackage[utf8]{inputenc}
\usepackage{lmodern}
\usepackage{amssymb}
\usepackage[hidelinks]{hyperref}
\usepackage{microtype}

\setlength{\parindent}{0pt}
\setlength{\parskip}{6pt}
\raggedbottom

\begin{document}

\begin{center}
    {\Large\bfseries Declaration of Generative AI Use}

    \vspace{8pt}
    {\bfseries Group 25: Reliable Multilingual Question Answering}

    \vspace{6pt}
    Haoyang Sun (mqb976)\\
    Kieran O'Neill (zcd890)\\
    Md Sultan Parvez (wvt774)

    \vspace{6pt}
    30 October 2026
\end{center}

\vspace{8pt}

$\boxtimes$ We have used generative AI as an aid/tool.

\paragraph{Tools and scope.}
Haoyang Sun and Md Sultan Parvez used
\href{https://chatgpt.com/}{ChatGPT (OpenAI)}.
Kieran O'neill used
\href{https://www.perplexity.ai/}{Perplexity}.
Specific underlying model versions are not consistently documented.
The following records distinguish individual uses; they do not imply
that every member used every tool or workflow.

\paragraph{Haoyang: purpose and work phases.}
GPT supported assignment interpretation and conceptual learning;
implementation suggestions, code generation/revision, code explanations
and debugging; supplementary-experiment design; numerical summaries,
tables and translations; and report organisation, prose drafting,
editing and LaTeX. It also supported group messages, progress/feedback
summaries and meeting action guides.

For Korean LM diagnostics, GPT translated selected questions into
English. For the translation extension, it translated Korean training
and validation questions; these translations were used without manual
editing or selection. For supplementary work, Haoyang independently
maintained an Excel analysis record documenting questions, experimental
settings, observed results, hypotheses and conclusions. It covers
punctuation sensitivity, transition-level LM evidence, the EM--F1 gap,
the translation comparison and challenge evaluation. Haoyang conducted
the experiments and inspected the results himself. GPT assisted with
supplementary-experiment design and implementation details, and helped
assess whether the evidence supported his proposed interpretations.

\paragraph{Haoyang: use and further work on outputs.}
Haoyang ran experiments locally, studied the implementations, inspected
examples, maintained the analysis record and proposed hypotheses and
interpretations. AI-generated code and prose were examined and revised
against implementation behaviour and experimental evidence. Discussion
with GPT helped refine the distinction between observations and causal
or semantic explanations. The gen-AI log has been exported for submission.

\paragraph{Parvez: purpose, phases and further use.}
ChatGPT assisted with Python code generation and improvement, debugging,
explanations of modelling/evaluation procedures, experimental-result
analysis, project planning and preliminary findings reports. This covered
BIO development, Korean language modelling, integrated-QA comparison,
question-conditioning ablation and challenge evaluation. Suggested code
was reviewed, adapted and tested locally; findings were checked against
saved predictions and actual outputs. AI-assisted reports supplied
working material for collaborative report preparation.

\paragraph{Kieran: purpose, phases and further use.}
Perplexity suggested the rule-based classifier, assisted debugging,
formulated challenge questions, and summarised individual work and
contributions for report preparation. Declared phases include Weeks 36,
43 and 44 and the contribution statement. Debugged code was adapted into
source code and executed in notebooks. Generated challenge questions were
used; translations were obtained with Google Translate. Kieran ran the
reported experiments and checked reported numbers against executed
outputs. His interaction logs are supplied as Markdown in
\texttt{genai-logs/zcd890} in the submission ZIP.

\end{document}